\input{preamble/preamble.tex}

\geometry{margin=0.85in}
\AtBeginDocument{\small}

\newcommand{\ExamNameField}{\noindent\textbf{Name:}\ \rule{0.7\linewidth}{0.4pt}\par\medskip}

\begin{document}
	\ExamTitleBlock{10th grade}{Term 2 C1 Catch-Up Activity (Solutions)}{}
	
	\ExamSection{Problems}
	\begin{ExamProblems}
		\item
		\subsection*{Problem description}
		A student entrepreneur must decide whether to rent a single pop-up kiosk for one weekend.
		Market demand may be \textit{high} or \textit{low}. The estimated profit (in hundreds of dollars)
		is 36 if demand is high and 10 if demand is low.
		
		\subsection*{Solution}
		\textbf{Step 1: Identify alternatives and states of nature.}
		The decision alternative is: Rent one kiosk.
		The states of nature are: High demand and Low demand.
		
		\textbf{Step 2: Construct the payoff table.}
		\begin{center}
			\textit{Payoff table}
			\begin{tabular}{l c c}
				\toprule
				State of nature & Rent kiosk \\
				\midrule
				High demand & 36 \\
				Low demand & 10 \\
				\bottomrule
			\end{tabular}
		\end{center}
		
		\textbf{Step 3: Explain payoff cells.}
		The value 36 means the profit when the kiosk is rented and demand is high.
		The value 10 means the profit when the kiosk is rented and demand is low.
		
		\textbf{Conclusion:} For C1, we have correctly identified one alternative, two states, and both payoff entries.
		
		\item
		\subsection*{Problem description}
		A coffee shop chain must choose one bean purchasing plan for the next month:
		Plan A (premium import) or Plan B (mixed sourcing).
		Consumer demand can be \textit{high} or \textit{low}. Profits are measured in thousands of dollars.
		For high demand, Plan A gives 84 and Plan B gives 72.
		For low demand, Plan A gives 22 and Plan B gives 34.
		
		\subsection*{Solution}
		\textbf{Step 1: Identify alternatives and states of nature.}
		Decision alternatives: Plan A, Plan B.
		States of nature: High demand, Low demand.
		
		\textbf{Step 2: Construct the payoff table.}
		\begin{center}
			\textit{Payoff table}
			\begin{tabular}{l c c}
				\toprule
				State of nature & Plan A & Plan B \\
				\midrule
				High demand & 84 & 72 \\
				Low demand & 22 & 34 \\
				\bottomrule
			\end{tabular}
		\end{center}
		
		\textbf{Step 3: Explain payoff cells.}
		Each cell is the profit for one plan under one demand state.
		For example, 72 is the profit of Plan B when demand is high.
		
		\textbf{Conclusion:} The table captures 2 alternatives, 2 states, and 4 clear payoff entries.
		
		\item
		\subsection*{Problem description}
		A small factory must choose a production mode for a seasonal product:
		Mode A (manual line) or Mode B (semi-automated line).
		The market may be \textit{strong}, \textit{stable}, or \textit{weak}.
		Estimated profits (in thousands of dollars) are:
		Mode A: 95, 60, 18; Mode B: 88, 66, 30.
		
		\subsection*{Solution}
		\textbf{Step 1: Identify alternatives and states of nature.}
		Decision alternatives: Mode A, Mode B.
		States of nature: Strong market, Stable market, Weak market.
		
		\textbf{Step 2: Construct the payoff table.}
		\begin{center}
			\textit{Payoff table}
			\begin{tabular}{l c c}
				\toprule
				State of nature & Mode A & Mode B \\
				\midrule
				Strong market & 95 & 88 \\
				Stable market & 60 & 66 \\
				Weak market & 18 & 30 \\
				\bottomrule
			\end{tabular}
		\end{center}
		
		\textbf{Step 3: Explain payoff cells.}
		Each entry gives profit for a specific mode under a specific market condition.
		For instance, 30 is Mode B profit when the market is weak.
		
		\textbf{Conclusion:} We correctly mapped 2 alternatives across 3 states with 6 payoff cells.
		
		\item
		\subsection*{Problem description}
		A delivery startup must choose one fleet strategy:
		Strategy A (owned vans), Strategy B (leased vans), or Strategy C (partner drivers).
		Fuel conditions for the quarter can be \textit{low}, \textit{medium}, or \textit{high} cost.
		Profits (in thousands of dollars) are:
		A: 120, 86, 44; B: 108, 92, 58; C: 96, 88, 70.
		
		\subsection*{Solution}
		\textbf{Step 1: Identify alternatives and states of nature.}
		Decision alternatives: Strategy A, Strategy B, Strategy C.
		States of nature: Low fuel cost, Medium fuel cost, High fuel cost.
		
		\textbf{Step 2: Construct the payoff table.}
		\begin{center}
			\textit{Payoff table}
			\begin{tabular}{l c c c}
				\toprule
				State of nature & Strategy A & Strategy B & Strategy C \\
				\midrule
				Low fuel cost & 120 & 108 & 96 \\
				Medium fuel cost & 86 & 92 & 88 \\
				High fuel cost & 44 & 58 & 70 \\
				\bottomrule
			\end{tabular}
		\end{center}
		
		\textbf{Step 3: Explain payoff cells.}
		A cell is the profit for one strategy under one fuel-cost state.
		For example, 58 means Strategy B under high fuel cost.
		
		\textbf{Conclusion:} The C1 structure is complete with 3 alternatives, 3 states, and 9 payoffs.
		
		\item
		\subsection*{Problem description}
		A supermarket must choose one inventory policy for imported fruit:
		Policy A (minimal stock), Policy B (balanced stock), or Policy C (buffer stock).
		Supply conditions can be \textit{very smooth}, \textit{smooth}, \textit{disrupted}, or \textit{severely disrupted}.
		Profits (in thousands of dollars) are:
		A: 74, 68, 36, 8; B: 82, 76, 50, 24; C: 86, 78, 58, 34.
		
		\subsection*{Solution}
		\textbf{Step 1: Identify alternatives and states of nature.}
		Decision alternatives: Policy A, Policy B, Policy C.
		States of nature: Very smooth, Smooth, Disrupted, Severely disrupted supply.
		
		\textbf{Step 2: Construct the payoff table.}
		\begin{center}
			\textit{Payoff table}
			\begin{tabular}{l c c c}
				\toprule
				State of nature & Policy A & Policy B & Policy C \\
				\midrule
				Very smooth supply & 74 & 82 & 86 \\
				Smooth supply & 68 & 76 & 78 \\
				Disrupted supply & 36 & 50 & 58 \\
				Severely disrupted supply & 8 & 24 & 34 \\
				\bottomrule
			\end{tabular}
		\end{center}
		
		\textbf{Step 3: Explain payoff cells.}
		Each payoff is the resulting profit when a specific inventory policy meets one supply condition.
		For example, 24 is Policy B profit under severely disrupted supply.
		
		\textbf{Conclusion:} The table correctly includes 3 alternatives, 4 states, and 12 payoff entries.
		
		\item
		\subsection*{Problem description}
		An electronics retailer must choose one launch format for a new device:
		Format A (online-only), Format B (store-only), Format C (hybrid), or Format D (limited rollout).
		The selling season can be \textit{boom}, \textit{good}, \textit{moderate}, or \textit{slow}.
		Profits (in thousands of dollars) are:
		A: 150, 110, 72, 30;
		B: 132, 118, 82, 48;
		C: 144, 124, 94, 62;
		D: 120, 108, 88, 66.
		
		\subsection*{Solution}
		\textbf{Step 1: Identify alternatives and states of nature.}
		Decision alternatives: Format A, Format B, Format C, Format D.
		States of nature: Boom, Good, Moderate, Slow season.
		
		\textbf{Step 2: Construct the payoff table.}
		\begin{center}
			\textit{Payoff table}
			\begin{tabular}{l c c c c}
				\toprule
				State of nature & Format A & Format B & Format C & Format D \\
				\midrule
				Boom season & 150 & 132 & 144 & 120 \\
				Good season & 110 & 118 & 124 & 108 \\
				Moderate season & 72 & 82 & 94 & 88 \\
				Slow season & 30 & 48 & 62 & 66 \\
				\bottomrule
			\end{tabular}
		\end{center}
		
		\textbf{Step 3: Explain payoff cells.}
		Every cell reports one format's profit in one season outcome.
		For example, 94 is the profit for Format C when the season is moderate.
		
		\textbf{Conclusion:} C1 is satisfied with 4 alternatives, 4 states, and 16 payoffs.
		
		\item
		\subsection*{Problem description}
		A grain exporter must choose one shipping contract:
		Contract A, Contract B, Contract C, or Contract D.
		Global freight conditions may be \textit{very favorable}, \textit{favorable}, \textit{neutral}, \textit{tight}, or \textit{very tight}.
		Profits (in thousands of dollars) are:
		A: 160, 138, 102, 66, 24;
		B: 152, 140, 114, 80, 40;
		C: 146, 136, 120, 92, 54;
		D: 132, 126, 116, 98, 70.
		
		\subsection*{Solution}
		\textbf{Step 1: Identify alternatives and states of nature.}
		Decision alternatives: Contract A, Contract B, Contract C, Contract D.
		States of nature: Very favorable, Favorable, Neutral, Tight, Very tight freight conditions.
		
		\textbf{Step 2: Construct the payoff table.}
		\begin{center}
			\textit{Payoff table}
			\begin{tabular}{l c c c c}
				\toprule
				State of nature & Contract A & Contract B & Contract C & Contract D \\
				\midrule
				Very favorable freight & 160 & 152 & 146 & 132 \\
				Favorable freight & 138 & 140 & 136 & 126 \\
				Neutral freight & 102 & 114 & 120 & 116 \\
				Tight freight & 66 & 80 & 92 & 98 \\
				Very tight freight & 24 & 40 & 54 & 70 \\
				\bottomrule
			\end{tabular}
		\end{center}
		
		\textbf{Step 3: Explain payoff cells.}
		Each entry shows the profit for one contract under one freight market condition.
		For instance, 70 is Contract D under very tight freight conditions.
		
		\textbf{Conclusion:} The payoff structure is complete with 4 alternatives, 5 states, and 20 cells.
		
		\item
		\subsection*{Problem description}
		A fashion retailer must choose one expansion plan:
		Plan A, Plan B, Plan C, Plan D, or Plan E.
		Urban demand can be \textit{very high}, \textit{high}, \textit{medium}, \textit{low}, or \textit{very low}.
		Profits (in thousands of dollars) are:
		A: 190, 150, 108, 58, 12;
		B: 182, 154, 120, 76, 34;
		C: 176, 152, 128, 90, 52;
		D: 168, 146, 126, 98, 64;
		E: 156, 140, 124, 104, 76.
		
		\subsection*{Solution}
		\textbf{Step 1: Identify alternatives and states of nature.}
		Decision alternatives: Plan A, Plan B, Plan C, Plan D, Plan E.
		States of nature: Very high, High, Medium, Low, Very low urban demand.
		
		\textbf{Step 2: Construct the payoff table.}
		\begin{center}
			\textit{Payoff table}
			\begin{tabular}{l c c c c c}
				\toprule
				State of nature & Plan A & Plan B & Plan C & Plan D & Plan E \\
				\midrule
				Very high demand & 190 & 182 & 176 & 168 & 156 \\
				High demand & 150 & 154 & 152 & 146 & 140 \\
				Medium demand & 108 & 120 & 128 & 126 & 124 \\
				Low demand & 58 & 76 & 90 & 98 & 104 \\
				Very low demand & 12 & 34 & 52 & 64 & 76 \\
				\bottomrule
			\end{tabular}
		\end{center}
		
		\textbf{Step 3: Explain payoff cells.}
		A payoff cell combines one expansion plan with one demand level.
		For example, 126 means Plan D when demand is medium.
		
		\textbf{Conclusion:} The C1 table lists 5 alternatives, 5 states, and 25 payoff values.
		
		\item
		\subsection*{Problem description}
		A regional energy distributor must choose one pricing package:
		Package A, Package B, Package C, Package D, or Package E.
		Weather-demand conditions can be \textit{extreme cold}, \textit{cold}, \textit{mild}, \textit{warm}, \textit{hot}, or \textit{extreme hot}.
		Profits (in thousands of dollars) are:
		A: 220, 184, 142, 96, 58, 22;
		B: 210, 186, 152, 110, 74, 40;
		C: 204, 182, 158, 122, 88, 54;
		D: 196, 176, 156, 128, 98, 68;
		E: 186, 170, 152, 130, 106, 82.
		
		\subsection*{Solution}
		\textbf{Step 1: Identify alternatives and states of nature.}
		Decision alternatives: Package A, Package B, Package C, Package D, Package E.
		States of nature: Extreme cold, Cold, Mild, Warm, Hot, Extreme hot.
		
		\textbf{Step 2: Construct the payoff table.}
		\begin{center}
			\textit{Payoff table}
			\begin{tabular}{l c c c c c}
				\toprule
				State of nature & Package A & Package B & Package C & Package D & Package E \\
				\midrule
				Extreme cold & 220 & 210 & 204 & 196 & 186 \\
				Cold & 184 & 186 & 182 & 176 & 170 \\
				Mild & 142 & 152 & 158 & 156 & 152 \\
				Warm & 96 & 110 & 122 & 128 & 130 \\
				Hot & 58 & 74 & 88 & 98 & 106 \\
				Extreme hot & 22 & 40 & 54 & 68 & 82 \\
				\bottomrule
			\end{tabular}
		\end{center}
		
		\textbf{Step 3: Explain payoff cells.}
		Each cell represents profit for one pricing package under one weather-demand state.
		For example, 88 is Package C under hot conditions.
		
		\textbf{Conclusion:} We have a complete C1 setup with 5 alternatives, 6 states, and 30 payoffs.
		
		\item
		\subsection*{Problem description}
		A national logistics group must choose one network design:
		Design A, Design B, Design C, Design D, Design E, or Design F.
		Macro-economic conditions can be \textit{rapid expansion}, \textit{steady growth}, \textit{flat market},
		\textit{mild contraction}, \textit{strong contraction}, or \textit{recovery transition}.
		Profits (in thousands of dollars) are:
		A: 260, 220, 170, 118, 70, 136;
		B: 248, 222, 182, 132, 86, 148;
		C: 242, 218, 188, 146, 102, 158;
		D: 234, 212, 190, 156, 118, 164;
		E: 226, 206, 188, 162, 130, 168;
		F: 214, 198, 184, 164, 142, 170.
		
		\subsection*{Solution}
		\textbf{Step 1: Identify alternatives and states of nature.}
		Decision alternatives: Design A, Design B, Design C, Design D, Design E, Design F.
		States of nature: Rapid expansion, Steady growth, Flat market, Mild contraction, Strong contraction, Recovery transition.
		
		\textbf{Step 2: Construct the payoff table.}
		\begin{center}
			\textit{Payoff table}
			\begin{tabular}{l c c c c c c}
				\toprule
				State of nature & Design A & Design B & Design C & Design D & Design E & Design F \\
				\midrule
				Rapid expansion & 260 & 248 & 242 & 234 & 226 & 214 \\
				Steady growth & 220 & 222 & 218 & 212 & 206 & 198 \\
				Flat market & 170 & 182 & 188 & 190 & 188 & 184 \\
				Mild contraction & 118 & 132 & 146 & 156 & 162 & 164 \\
				Strong contraction & 70 & 86 & 102 & 118 & 130 & 142 \\
				Recovery transition & 136 & 148 & 158 & 164 & 168 & 170 \\
				\bottomrule
			\end{tabular}
		\end{center}
		
		\textbf{Step 3: Explain payoff cells.}
		Each of the 36 cells gives the profit for one design in one macro-economic state.
		For instance, 164 is the payoff for Design D during recovery transition.
		
		\textbf{Conclusion:} The final problem satisfies C1 with 6 alternatives, 6 states, and 36 coherent payoff entries.
	\end{ExamProblems}
\end{document}
